\section{Finite Frozen-Pool ERM}
\label{sec:model-selection}

Fix an architecture $\arch$ together with measurement budget $M$, pool size
$S$, and readout radius $\Lambda$.  Each candidate runs the same declared
architecture but with an independently frozen realization.  Selection is over
this finite pool only: no gate axis, angle, or contraction rate is optimized.

\subsection{Frozen pool and cached measurements}
\label{sec:frozen-pool-fitting}

Draw
\[
    \design_1,\ldots,\design_S
    \overset{\mathrm{i.i.d.}}{\sim}\designLaw_{\arch},
\]
independently of the labeled sample and freeze the pool.  For every observed
window and every candidate $s$, generate
$\measuredFeatureMap_{\design_s}$ once and cache it.  With the cycle cover,
the complete measured training pool costs $9NMSR$ independent state
preparations/branch-circuit runs under the reusable-$n$-qubit convention; all
subsequent readout fits and design selection are classical.  As clarified in
\cref{sec:topology-observables}, this acquisition count is separate from the
$O(w)$ logical depth and $O(n)$ width of one trajectory.

\subsection{Joint empirical-risk selection}
\label{sec:full-sample-erm}
\label{sec:structural-risk-selection}

QuaRK selects the frozen design and its readout by one finite-pool empirical
risk minimization,
\begin{equation}
    (\widehat s,\selectedReadout)
    \in\argmin_{\substack{1\leq s\leq S\\
                          h\in\RKHS_{m_{\arch}}^\Lambda}}
       \erisk_{\design_s,w}(h),
    \qquad
    \selectedDesign:=\design_{\widehat s},
    \label{eq:frozen-pool-erm}
\end{equation}
with a deterministic candidate tie rule and the minimum-RKHS-norm minimizer
within a tied candidate.  Equivalently, one fits an Ivanov-regularized kernel
least-squares readout for each frozen design and keeps the smallest measured
empirical risk.  The representer theorem reduces every fit to a finite convex
problem \citep{PageGrunewalder2019,scholkopf2001generalized}.

Using the same sample for fitting and selection is deliberate rather than an
assumption that overfitting cannot occur.  The uniform deviation result in
\cref{sec:dependent-generalization} holds simultaneously over all $S$ designs
and the entire Ivanov ball.  It therefore controls this data-dependent ERM
directly: searching more frozen designs costs only a $\log S$ term, while
readout adaptivity is controlled by $\Lambda$.  This makes the statistical
price of frozen-pool search explicit without introducing a second validation
procedure.

The returned predictor is
$\selectedReadout\circ\measuredFeatureMap_{\selectedDesign}$ operationally,
with $\selectedReadout\circ\featureMap_{\selectedDesign,\infty}$ as its causal
exact-feature counterpart.  The role of $S$ is quantified in
\cref{sec:pool-coverage}: it increases the chance that the frozen pool contains
a near-best representation, while the same increase is charged statistically
through the uniform $\log S$ term.
